
\documentclass{article}
\usepackage{colortbl}
\usepackage{makecell}
\usepackage{multirow}
\usepackage{supertabular}

\begin{document}

\newcounter{utterance}

\centering \large Interaction Transcript for game `cladder', experiment `full\_v1.5\_default', episode 6572 with Qwen3.5{-}9B.
\vspace{24pt}

{ \footnotesize  \setcounter{utterance}{1}
\setlength{\tabcolsep}{0pt}
\begin{supertabular}{c@{$\;$}|p{.15\linewidth}@{}p{.15\linewidth}p{.15\linewidth}p{.15\linewidth}p{.15\linewidth}p{.15\linewidth}}
    \# & \multicolumn{2}{c}{Player} && \multicolumn{2}{c}{Game Master} \\
    \hline

    \theutterance \stepcounter{utterance}  
    & & & \multicolumn{4}{p{0.6\linewidth}}{
        \cellcolor[rgb]{0.9,0.9,0.9}{
            \makecell[{{p{\linewidth}}}]{
                \texttt{\tiny{[P1$\langle$GM]}}
                \texttt{You are an expert in causal inference. The following question is not a typical commonsense query, but rather a meticulously designed question created by a professor specializing in causal inference, intended to assess the students' mastery of the course content.} \\
\texttt{Imagine a self{-}contained, hypothetical world with only the following conditions, and without any unmentioned factors or causal relationships: Kwox has a direct effect on kwoz. Swoq has a direct effect on kwoz. For those who are not kwox, the probability of kwoz is 82\%. For those who are kwox, the probability of kwoz is 32\%. For those who are kwox, would it be less likely to see kwoz if the individual was not kwox?} \\
\texttt{Start your answer with 'Yes' or 'No', followed by additional reasoning or evidence to support your explanation.} \\
            }
        }
    }
     \\ \\

    \theutterance \stepcounter{utterance}  
    & \multicolumn{4}{p{0.6\linewidth}}{
        \cellcolor[rgb]{0.561,0.737,0.561}{
            \makecell[{{p{\linewidth}}}]{
                \texttt{\tiny{[P1$\rangle$GM]}}
                \texttt{No} \\
\\ 
\texttt{**Reasoning:**} \\
\\ 
\texttt{To determine the answer, we must compare the probability of the outcome \$kwoz\$ under two specific conditions for an individual who is \$kwox\$:} \\
\texttt{1.  **The actual condition:** The individual is \$kwox\$.} \\
\texttt{2.  **The counterfactual condition:** The individual is *not* \$kwox\$ (while holding their other characteristics constant, specifically their status as \$kwox\$ in the real world, though the prompt implies a comparison of the marginal probabilities defined by the variable \$kwox\$).} \\
\\ 
\texttt{Based on the data provided in the hypothetical world:} \\
\texttt{*   \$P(kwoz \textbackslash{}mid kwox) = 32\textbackslash{}\%\$} \\
\texttt{*   \$P(kwoz \textbackslash{}mid \textbackslash{}neg kwox) = 82\textbackslash{}\%\$} \\
\\ 
\texttt{The question asks: "For those who are kwox, would it be less likely to see kwoz if the individual was not kwox?"} \\
\\ 
\texttt{This is asking if \$P(kwoz \textbackslash{}mid \textbackslash{}neg kwox) < P(kwoz \textbackslash{}mid kwox)\$.} \\
\texttt{Substituting the values: Is \$82\textbackslash{}\% < 32\textbackslash{}\%\$?} \\
\\ 
\texttt{This statement is false. In fact, the probability of seeing \$kwoz\$ is significantly **higher** (82\%) if the individual were not \$kwox\$, compared to when they are \$kwox\$ (32\%). Therefore, it would be *more* likely, not less likely, to see \$kwoz\$ in the counterfactual scenario where the individual is not \$kwox\$.} \\
\\ 
\texttt{The direct effect of \$kwox\$ on \$kwoz\$ is negative (since the probability drops from 82\% to 32\% when \$kwox\$ is present), but the question asks about the likelihood of the outcome in the absence of the cause relative to the presence of the cause. Since the absence of the cause yields a higher probability, the answer is No.} \\
            }
        }
    }
    & & \\ \\

    \theutterance \stepcounter{utterance}  
    & & & \multicolumn{2}{p{0.3\linewidth}}{
        \cellcolor[rgb]{0.9,0.9,0.9}{
            \makecell[{{p{\linewidth}}}]{
                \texttt{\tiny{[GM$|$GM]}}
                \texttt{No} \\
            }
        }
    }
    & & \\ \\

    \theutterance \stepcounter{utterance}  
    & & & \multicolumn{2}{p{0.3\linewidth}}{
        \cellcolor[rgb]{0.9,0.9,0.9}{
            \makecell[{{p{\linewidth}}}]{
                \texttt{\tiny{[GM$|$GM]}}
                \texttt{no} \\
            }
        }
    }
    & & \\ \\

    \theutterance \stepcounter{utterance}  
    & & & \multicolumn{2}{p{0.3\linewidth}}{
        \cellcolor[rgb]{0.9,0.9,0.9}{
            \makecell[{{p{\linewidth}}}]{
                \texttt{\tiny{[GM$|$GM]}}
                \texttt{game\_result = WIN} \\
            }
        }
    }
    & & \\ \\

\end{supertabular}
}

\end{document}
